% Generated draft: requires mathematical and source review.
\documentclass[11pt]{article}
\usepackage[margin=25mm]{geometry}
\usepackage{fontspec}
\setmainfont{Latin Modern Roman}
\usepackage{amsmath,amssymb,mathtools,unicode-math}
\setmathfont{Latin Modern Math}
\usepackage{enumitem,needspace,xurl,hyperref}
\setlength{\parindent}{0pt}
\setlength{\parskip}{5pt}
\setlength{\emergencystretch}{3em}
\newcommand{\N}{\mathbb{N}}
\newcommand{\Z}{\mathbb{Z}}
\newcommand{\Q}{\mathbb{Q}}
\newcommand{\R}{\mathbb{R}}
\newcommand{\C}{\mathbb{C}}
\newcommand{\F}{\mathbb{F}}
\newcommand{\A}{\mathbb{A}}
\newcommand{\PP}{\mathbb{P}}
\newcommand{\E}{\mathbb{E}}
\newcommand{\eps}{\varepsilon}
\newcommand{\dd}{\,\mathrm{d}}
\begin{document}
\section*{UnsolvedMath ID: 30000483}
\textbf{Source problem code:} catalogue entry
\subsection{Definitions and mathematical statement}
Let \(H\) be a Noetherian AS-Gorenstein Hopf algebra over the source's base field \(k\), with coproduct \(\Delta\), counit \(\epsilon\), and antipode \(S\). In the usual AS-Gorenstein convention, the left and right injective dimensions are finite and equal to \(d\), with the corresponding \(\operatorname{Ext}^i(k,H)\) vanishing outside degree \(d\) and one-dimensional in that degree. The trivial modules use the counit. The record does not reproduce the full convention or identify which one-dimensional homological-integral module defines its character \(\pi:H\to k\).

For a specified character \(\pi\), standard winding-automorphism conventions are
\[
\tau_\pi^\ell(h)=\sum\pi(h_{(1)})h_{(2)},\qquad
\tau_\pi^r(h)=\sum h_{(1)}\pi(h_{(2)}),
\quad\Delta(h)=\sum h_{(1)}\otimes h_{(2)}.
\]
The side and inverse conventions must be matched to the source. Its supplied fourth-power formula is
\[
S^4=\gamma\circ\tau_\pi^r\circ(\tau_\pi^\ell)^{-1},
\qquad\gamma\in\operatorname{Inn}(H).
\]
Here \(S^4\) is fourfold composition and an inner automorphism has the form \(\gamma(h)=uhu^{-1}\) for a unit \(u\in H^\times\). The problem is to determine \(\gamma\) explicitly in general, equivalently to give an explicit implementing unit or structural description satisfying
\[
S^4(h)=u\bigl(\tau_\pi^r((\tau_\pi^\ell)^{-1}(h))\bigr)u^{-1}\qquad(h\in H).
\]
The implementing unit is only determined up to multiplication by a central unit. With \(S\) and \(\pi\) fixed, the displayed identity already fixes \(\gamma\) formally; merely solving the equation for a composite of known automorphisms does not explain the intended request for an explicit inner description.

The supplied literature assessment checked 2026-08-21 reports existence of an inner automorphism in the formula but says the cited results do not give a general explicit determination. Thus existence of some \(u\) is source background, not the requested new conclusion. The desired form of explicit determination, and the coefficient and integral conventions on which it depends, remain to be recovered.
\subsection{Short English statement}
Give a structural or explicit inner description of the correction in the fourth-power antipode formula, beyond the supplied existence theorem. The character and winding conventions matter, and an implementing unit is only defined up to central units.
\subsection{Sources}
\begin{itemize}
\item Local UnsolvedMath registry, permanent ID 30000483, catalogue entry, \emph{Inner Automorphisms in the Fourth-Power Antipode Formula}.
\item \emph{Interactions between Algebraic Geometry and Noncommutative Algebra} (2007), Oberwolfach Reports, pp. 1319--1384, DOI \url{https://doi.org/10.4171/owr/2006/23}.
\item Lv--Oh--Wang--Yu, supplied label ``generalized S-four formula'': \url{https://arxiv.org/abs/1704.05731}.
\item Supplied statement-verification claim and literature triage checked 2026-08-21 distinguish the existence formula from explicit determination.
\end{itemize}
\end{document}
